\chapter{Modèle de données et persistance}
\label{ch:donnees}

\section{Entités principales}

\begin{description}
  \item[\texttt{Asset}] Représente un instrument : symbole (unique), nom, type (\texttt{Stock} / \texttt{Crypto}), prix courant. Collections de navigation vers les transactions et l'historique de prix.
  \item[\texttt{Transaction}] Mouvement d'achat ou de vente : quantité, prix unitaire, date, montant total dérivé ; clé étrangère vers \texttt{Asset}.
  \item[\texttt{PriceHistory}] Série temporelle de clôtures simulées ou issues d'une évolution de marché ; contrainte d'unicité \((\texttt{AssetId}, \texttt{Date})\).
  \item[\texttt{PortfolioBudget}] Budget initial et liquidité disponible (mise à jour lors des opérations).
  \item[\texttt{PortfolioValuePoint}] Valeur totale du portefeuille (cash + positions au cours du jour) pour une date donnée ; index unique sur la date.
\end{description}

\section{Relations et intégrité}

Les relations \texttt{Asset} $\rightarrow$ \texttt{Transaction} et \texttt{Asset} $\rightarrow$ \texttt{PriceHistory} sont en \textbf{one-to-many} avec suppression en cascade côté dépendant, configurées dans \texttt{OnModelCreating}. L'index unique sur \texttt{Asset.Symbol} garantit l'absence de doublons de tickers.

\section{Migrations}

Les migrations EF Core sont stockées dans \texttt{GestionPortefeuille.Infrastructure/Migrations/}, notamment :
\begin{itemize}[leftmargin=*]
  \item création initiale des tables (\texttt{InitialCreate}) ;
  \item ajout de \texttt{PortfolioValuePoints} (\texttt{AddPortfolioValuePoint}).
\end{itemize}

Au démarrage de l'application, \texttt{Database.Migrate()} applique automatiquement les migrations en attente, ce qui facilite la démonstration sur une machine neuve. Un \emph{seed} insère un budget par défaut et des actifs de démonstration si la base est vide.

\section{Choix SQLite}

SQLite convient à un projet étudiant : fichier unique, pas de serveur à installer, sauvegarde du dépôt possible en excluant \texttt{*.db} via \texttt{.gitignore}. Pour une mise en production, une base serveur (PostgreSQL, SQL Server) serait préférable (concurrence, sauvegardes, taille).
